\section{弯折拐角壁面相似解建模}
讨论二维流场两个互相垂直壁面的边界层相似解的求取问题。

方便起见，在这个问题中建立两个坐标系，第一个坐标系$xy$坐标系描述竖直方向壁面的流动情况，第二个坐标系$x'y'$，描述水平方向壁面的流动情况。两个坐标系坐标之间的关系为$x=x',y=H-y'$。

对两个坐标系分别建立二维稳态黏性层流NS方程。

\subsection{坐标系关系说明}
\begin{itemize}
    \item 水平壁面坐标系$(x,y)$:
    \begin{itemize}
        \item $x$：水平坐标，平行于地面；
        \item $y$：竖直坐标，沿着竖直壁面向上；
    \end{itemize}
    \item 竖直壁面坐标系$(x',y')$:
    \begin{itemize}
        \item $x'$：水平坐标，与x相同；
        \item $y'$：竖直坐标，沿着竖直壁面向下；
    \end{itemize}
\end{itemize}
它们之间的关系为：
\begin{align*}
    x=x',y=H-y'
\end{align*}
$H$为竖壁长度以及两个坐标原点竖向的距离。
微分变换为：
\begin{align*}
    \frac{\partial}{\partial x}=\frac{\partial }{\partial x'},\frac{\partial}{\partial y}=-\frac{\partial}{\partial y'}
\end{align*}

\subsection{在两个坐标系下建立二维稳态不可压黏性层流NS方程}
$(x,y)$坐标系下，设速度分量分别为$u(x,y),v(x,y)$，分别为$x,y$方向的分速度。
\begin{align}
    \frac{\partial u}{\partial x}+\frac{\partial v}{\partial y}&=0\\
    u\frac{\partial u}{\partial x}+v\frac{\partial u}{\partial y}&=-\frac{1}{\rho}\frac{\partial p}{\partial x}+\nu\left(\frac{\partial^{2} u}{\partial x^{2}}+\frac{\partial^{2} u}{\partial y^{2}}\right)\\
    u\frac{\partial v}{\partial x}+v\frac{\partial v}{\partial y}&=-\frac{1}{\rho}\frac{\partial p}{\partial y}+\nu\left(\frac{\partial^{2} v}{\partial x^{2}}+\frac{\partial^{2} v}{\partial y^{2}}\right)
\end{align}

同理，$(x',y')$坐标系下，设速度分量分别为$u'(x',y'),v'(x',y')$，分别为$x',y'$方向的分速度。

根据第一节的坐标关系，有如下变换：
\begin{align*}
    u'(x',y')&=u(x,y),&v'(x',y')=-v(x,y)\\
    \frac{\partial u}{\partial x}&=\frac{\partial u'}{\partial x'}, &\frac{\partial u}{\partial y}=-\frac{\partial u'}{\partial y'}\\
    \frac{\partial v}{\partial x}&=-\frac{\partial v'}{\partial x'}, &\frac{\partial v}{\partial y}=\frac{\partial v'}{\partial y'}\\
    \frac{\partial^{2} u}{\partial x^{2}}&=\frac{\partial^{2} u'}{\partial x'^{2}}, &\frac{\partial^{2} u}{\partial y^{2}}=\frac{\partial^{2} u'}{\partial y'^{2}}\\
    \frac{\partial^{2} v}{\partial x^{2}}&=\frac{\partial^{2} (-v')}{\partial x'^{2}}, & \frac{\partial^{2} v}{\partial y^{2}}=-\frac{\partial^{2} v'}{\partial y'^{2}}\\
    \frac{\partial p}{\partial x}&=\frac{\partial p'}{\partial x'}, &\frac{\partial p}{\partial y}=-\frac{\partial p'}{\partial y'}
\end{align*}

代入$(x,y)$坐标系下的NS方程即得到$(x',y')$坐标系下的NS方程
\begin{align}
    \frac{\partial u'}{\partial x'}+\frac{\partial v'}{\partial y'}&=0\\
    u'\frac{\partial u'}{\partial x'}+v'\frac{\partial u'}{\partial y'}&=-\frac{1}{\rho}\frac{\partial p'}{\partial x'}+\nu\left(\frac{\partial^{2} u'}{\partial x'^{2}}+\frac{\partial^{2} u'}{\partial y'^{2}}\right)\\
    u'\frac{\partial v'}{\partial x'}+v'\frac{\partial v'}{\partial y'}&=-\frac{1}{\rho}\frac{\partial p'}{\partial y'}+\nu\left(\frac{\partial^{2} v'}{\partial x'^{2}}+\frac{\partial^{2} v'}{\partial y'^{2}}\right)
\end{align}
与$(x,y)$坐标系下的形式完全一样，形式保持不变。

\subsection{对竖直壁面的流动简化到壁射流边界层方程}
为了方便起见，都使用$x,y$作为坐标符号。

引入流函数，则连续方程自动满足，动量方程表示为：
\begin{align}
    \frac{\partial \psi}{\partial y}\frac{\partial^{2} \psi}{\partial x\partial y}-\frac{\partial \psi}{\partial x}\frac{\partial^{2} \psi}{\partial y^{2}}&=-\frac{1}{\rho}\frac{\partial p}{\partial x}+\nu(\frac{\partial^{3} \psi}{\partial x^{2}\partial y}+\frac{\partial^{3} \psi}{\partial y^{3}})\\
    \frac{\partial \psi}{\partial y}\frac{\partial^{2} \psi}{\partial x^{2}}-\frac{\partial \psi}{\partial x}\frac{\partial^{2} \psi}{\partial x\partial y}&=\frac{1}{\rho}\frac{\partial p}{\partial y}+\nu(\frac{\partial^{3} \psi}{\partial x^{3}}+\frac{\partial^{3} \psi}{\partial x\partial y^{2}})
\end{align}

$x,y$分别表示壁面法向与壁面切向，简化后，$y$方向动量方程为：
\begin{align}
    \frac{\partial \psi}{\partial y}\frac{\partial^{2} \psi}{\partial x^{2}}-\frac{\partial \psi}{\partial x}\frac{\partial^{2} \psi}{\partial x\partial y}&=\frac{1}{\rho}\frac{\partial p}{\partial y}+\nu\frac{\partial^{3} \psi}{\partial x^{3}}\\
    0&=\frac{\partial p}{\partial x}
\end{align}
结论，简化后的边界层方程中，$p$是切向坐标的单值函数。

设局部特征尺度为$U(y),L(y)$，定义相似解的形式为：
\begin{align}
    \psi=ULf(\frac{x}{L})=ULf(\xi)
\end{align}

相关变量罗列：
\begin{align}
    \frac{\partial \psi}{\partial x}&=Uf'\\
    \frac{\partial \psi}{\partial y}&=f\left(\frac{\mathrm{d}U}{\mathrm{d}x}L+\frac{\mathrm{d}L}{\mathrm{d}y}U\right)-\frac{Ux}{L}\frac{\mathrm{d}L}{\mathrm{d}y}f'\\
    \frac{\partial^{2} \psi}{\partial x^{2}}&=\frac{U}{L}f''\\
  \frac{\partial^{3} \psi}{\partial x^{3}}&=\frac{U}{L^{2}}f'''\\
      \frac{\partial^{2} \psi}{\partial x\partial y}&=f'\frac{\mathrm{d}U}{\mathrm{d}y}-\frac{Ux}{L^{2}}\frac{\mathrm{d}L}{\mathrm{d}y}f'  
\end{align}

代入并化简合并得到ODE为：
\begin{align}
    f'''-ff''\left(\frac{L^{2}}{\nu}\frac{\mathrm{d}U}{\mathrm{d}y}+\frac{UL}{\nu}\frac{\mathrm{d}L}{\mathrm{d}y}\right)+f'f'\frac{L^{2}}{\nu}\frac{\mathrm{d}U}{\mathrm{d}y}+\frac{L^{2}}{U\mu}\frac{\partial p}{\partial y}=0
\end{align}
若要使相似解存在，则需要满足以下条件(自相似解存在的充分条件):
\begin{align}
    \frac{UL}{\nu}\frac{\mathrm{d}L}{\mathrm{d}y}=\alpha\\
    \frac{L^{2}}{\nu}\frac{\mathrm{d}U}{\mathrm{d}y}=\beta\\
    \frac{L^{2}}{U\mu}\frac{\mathrm{d} p}{\mathrm{d} y}=\gamma
\end{align}
其中，$\alpha,\beta,\gamma$为常数。

则原ODE变为
\begin{align}
    \label{eq-ODE}
    f'''-(\beta+\alpha)ff''+\beta f'f'+\gamma=0
\end{align}

将条件方程组中（2）（3）时相除得到：
\begin{align}
    p(U)=\frac{\gamma \rho}{2\beta}U^{2}+p_{0}
\end{align}

上式表示，在黏性主导流动中，为保持自相似结构所要求的压强-速度协调条件，数学形式类比于 Bernoulli 方程。

\subsubsection{幂律假设}
若假设局部特征尺度有以下形式：
\begin{align}
    L(y)=Ay^{m},U(y)=By^{n}
\end{align}

代入条件方程组得到：
\begin{align}
    L(y)&=Ay^{\frac{\alpha}{\beta+2\alpha}}\\
    U(y)&=\frac{\nu(\beta+2\alpha)}{A^{2}}y^{\frac{\beta}{\beta+2\alpha}}\\
    p(y)&=\gamma\frac{\nu(\beta+2\alpha)}{A^{4}}y^{\frac{\beta-2\alpha}{\beta+2\alpha}}
\end{align}

整理后得到：
\begin{align}
    p=\gamma \frac{U}{L^{2}}
\end{align}

\subsubsection{指数律假设}
若假设局部尺度特征有以下形式：
\begin{align}
    L(y)=Ae^{my},U(y)=Be^{ny}
\end{align}

代入条件方程组得：
\begin{align}
    L(y)&=Ae^{\frac{\alpha n}{\beta}y}\\
    U(y)&=Be^{\frac{-2\alpha n}{\beta}y}\\
    p(y)&=\frac{\gamma \mu B}{2nA^{2}}e^{2ny}+p_{0}
\end{align}

\subsection{自由射流（有压力梯度）}
现在分析有压力梯度影响的自由射流问题的守恒量形式。

边界层控制方程：
\begin{align}
    \frac{\partial u}{\partial x}+\frac{\partial v}{\partial y}&=0\\
    \rho\left(u\frac{\partial u}{\partial x}+v\frac{\partial u}{\partial y}\right)&=-\frac{\partial p}{\partial x}+\mu\frac{\partial^{2} u}{\partial y^{2}}
\end{align}
边界条件
\begin{align}
    y&=0, 0<x<H: u=0,v=0,\psi=0\\
    y&=\pm\infty,0<x<H: u=0, \frac{\partial u}{\partial y}=0
\end{align}
对连续方程变形处理后代入NS方程：
\begin{align}
    \rho u \frac{\partial u}{\partial x}+\rho u\frac{\partial v}{\partial y}&=0\\
    \rho\left(\frac{\partial uu}{\partial x}+\frac{\partial uv}{\partial y}\right)&=-\frac{\partial p}{\partial x}+\mu\frac{\partial \tau}{\partial y}
\end{align}

对变形后动量方程沿射流截面从$-\infty$到$\infty$积分：
\begin{align}
    \rho\left(\int_{-\infty}^{\infty}\frac{\partial uu}{\partial x}\mathrm{d}y+\int_{-\infty}^{\infty}\frac{\partial uv}{\partial y}\mathrm{d}y\right)&=-\int_{-\infty}^{\infty}\frac{\partial p}{\partial x}\mathrm{d}y+\int_{-\infty}^{\infty}\mu\frac{\partial \tau}{\partial y}\mathrm{d}y
\end{align}
其中等式左右两边的第二项可以根据边界条件得出均为0，则简化后的方程为
\begin{align}
    \rho\frac{\mathrm{d}}{\mathrm{d}x}\int_{-\infty}^{\infty}(uu+p)\mathrm{d}y=0\\
    \rho\int_{-\infty}^{\infty}(uu+p)\mathrm{d}y=\cst=J^{*}
\end{align}
若忽略x方向上的压力变换，则压力梯度项为零，上式退化为：
\begin{align}
     \rho\int_{-\infty}^{\infty}uu\mathrm{d}y=\cst=J
\end{align}
此为无压力梯度自由射流动量通量守恒表达式。

\subsection{自相似解代入守恒量}

自相似变量表示为：
\begin{align}
    u(x,y)&=U(x)f'(\eta)\\
    \eta&=\frac{y}{L(x)}\\
    \mathrm{d}y&=L(x)\mathrm{d}\eta
\end{align}
代入守恒量定义式
\begin{align}
    \rho\int_{-\infty}^{\infty}(U^{2}f'^{2}+p)L\mathrm{d}\eta=\cst\\
    \rho U^{2}L\int_{-\infty}^{\infty}f'^{2}\mathrm{d}\eta+\rho L\int_{-\infty}^{\infty}p\mathrm{d}\eta=\cst
\end{align}
